\subsection{Congruent real projections at low free rank}
\label{sec:lr-expanded}

We now specify the finite verification used when the original signed
matching presentation has free rank one or two. All lattices in this
subsection are integral lattices in the same labelled $\Z^{10}$.
Only the auxiliary height arrangement is a rational arrangement.

\begin{lemma}[LR]\label{lem:lr}
Let $M$ be a full signed matching matrix and let
$d=10-\rank_{\mathbb Q}M\in\{1,2\}$. Suppose that one fixed relabelling,
reflection and reanchoring makes every real solution satisfy labelled
$B=A$. Every cyclic realization having six distinct points on each side
and different $\TI$ classes is a specialization of a Bloom, D, L2,
L3*, L4, L5 or L7 construction.
\end{lemma}

The proof consists of the mathematical reductions below and the stated
finite acceptance predicates. The counts describe accepted existing
certificate packages; they are not assumptions about a numerical trend.

\paragraph{The height direction arrangement.}
Put $A=(0,a_1,\ldots,a_5)$ after the fixed normalization of $B$.
Before anchoring, let $\alpha_{ij}=e_j-e_i\in\Z^6$, $i\ne j$.
Every substituted matching row is a difference of two such oriented
edges. Its nonzero primitive direction, considered up to sign, is a
permutation of
\[
 (1,-1,0,0,0,0),\qquad (2,-1,-1,0,0,0),\qquad
 (1,1,-1,-1,0,0).
\]
There are respectively $15$, $60$, and $45$ directions. Denote their
union by $\mathcal D$. Let $S$ be the rational span of the substituted
rows. The space of their height solutions is $S^\perp$, modulo the
constant vector. The projection from the original real solution space
to the anchored $A$ coordinates is an isomorphism: it is injective
because $B=A$, and substitution gives its inverse. Consequently
$\dim S=5-d$, hence $\dim S$ is three or four. Dividing out integer
content here describes $S$; it does not change any matching lattice.

\begin{lemma}[Exhaustive height representatives]
Every rank-$r$ span generated by $\mathcal D$ occurs, up to a vertex
permutation, by adjoining a direction to a representative of rank $r-1$.
For every such span $S$ there is an integral vector $h$, with $h_0=0$,
such that $v\cdot h=0$ for $v\in\mathcal D$ exactly when $v\in S$.
\end{lemma}
\begin{proof}
Choose a basis of $S$ from its generating directions and remove its last
member. A permutation taking the remaining span to a representative
also takes the removed member into $\mathcal D$, since $\mathcal D$ is
permutation invariant. Induction proves the first assertion, starting
at the zero span. For the second, each $v\notin S$ defines a proper
rational hyperplane in $S^\perp$. A finite union of proper hyperplanes
cannot cover this vector space. Choose a rational vector outside their
union, subtract its first coordinate from every coordinate, and clear
denominators. All directions have coordinate sum zero, so anchoring
does not alter the required evaluations.
\end{proof}

The finite height check starts with all spans of ranks zero, one and
two, and verifies the extension induction by exact rational arithmetic.
For each saved representative it applies all $720$ vertex permutations,
checks orbit size and disjointness, and checks every extension of every
preceding representative. Every saved orbit must meet that extension
list. Independently, a span can be represented by its primitive exterior
minors: two full-rank row lists represent the same rational span exactly
when their exterior coordinate vectors are proportional. The accepted
rank-three table has $43\,770$ spans in $104$ orbits; the rank-four table
has $116\,401$ spans in $211$ orbits. Each saved $h$ is checked against
all $120$ directions using exact dot products and exact span membership.
Thus equal heights and equal signed edge heights are precisely the
coincidences forced by the representative span, without sampling error.

\paragraph{Roots retain the original integral relations.}
For one fixed $h$, orient every nonzero-height edge upward and group
these edges by their positive height differences. If a bin has $m$
edges, its compatible bijections are the $m!$ permutations. A full
matching respecting $h$ restricts to one element of the product of
these bin-permutation sets. Equal-height edges are left for completion.
Relabellings inside each equal-height fibre act independently on the
source and target lists. They preserve the completion problem and
are integral invertible changes of anchored generators. In particular,
when the zero-labelled vertex moves, reanchoring subtracts its new
coordinate; the inverse relabelling gives an integral inverse.

For an upward source edge $(i,j)$ and target edge $(k,l)$ the root row
is $\alpha_{ij}^A-\alpha_{kl}^B$ in anchored coordinates. Upward order
need not have $i<j$. To use the fixed $i<j$ convention, multiply the
whole row by $-1$ if the source edge is reversed and transfer that
sign to the target; reverse the target edge as necessary and change
its sign accordingly. These operations leave its integral row lattice
unchanged. No equation $v\cdot h=0$ is added simply because it holds.

The root acceptance check reconstructs all bins and all within-fibre
vertex permutations. For each saved bin permutation $p$, it constructs
its complete orbit $\{\beta p\alpha\}$ under the independent actions,
checks the recorded orbit size, and checks disjointness from earlier
orbits. Their union must have cardinality $\prod m!$. It reconstructs
every saved root matrix from its edge permutation, rather than trusting
the saved matrix. Across the $315$ height strata this accounts for
$4\,822\,240$ cross-edge bijections in $328\,473$ root orbits.

\paragraph{Exact group coordinates and cache equality.}
A node has a row matrix $H\in\Z^{k\times10}$ and represents
$G_H=\Z^{10}/\row_\Z(H)$. Its Smith certificate consists of
$U\in\Z^{k\times k}$ and $V\in\Z^{10\times10}$ with
\[
 \det U,\det V\in\{1,-1\},\qquad UHV=D.
\]
Here $D$ is a rectangular diagonal matrix with $r$ nonzero entries,
where $r=\rank H$. Thus $U$ is not necessarily $15\times15$:
its dimension changes with the node's actual row count. Keeping the
coordinates of rows of $V$ at indices $i<r$ with $|D_{ii}|>1$, reduced
modulo $|D_{ii}|$, and all indices $i\ge r$ gives the ten generator
coordinates in $T\oplus\Z^{10-r}$. Entries $|D_{ii}|=1$ contribute
no coordinate. The two anchors are zero. The verifier checks the
matrix identity, dimensions, determinants, rank, diagonal data and all
twelve recorded point coordinates by integer arithmetic.

The producer caches the row-Hermite matrix of the relation lattice.
This cache is local to one fixed height stratum and the same labelled
generators; neither rational span equality nor an unrecorded abstract
group isomorphism permits a merge. For every root normalization and
every parent-plus-branch normalization, the package supplies integral
matrices $C,C'$ satisfying
\[
 C H_{\rm source}=H_{\rm target},\qquad
 C'H_{\rm target}=H_{\rm source}.
\]
These equalities prove equality of integral row lattices in both
directions, including torsion, without trusting a Hermite routine.
Canonical minimality is useful for efficiency, not for the coverage
proof: two-way lattice equality is the required mathematical contract.

\paragraph{Paired cancellation and exhaustive branching.}
In $G_H$ form a bin for each unsigned class $\{z,-z\}$ of the fifteen
edge differences, retaining its labelled edge occurrences. Give the
class the deterministic representative $\min(z,-z)$ in Smith coordinates.
Cancel the minimum of the two multiplicities for each common class.
This cancels equally many occurrences on the two sides, including
order-two classes; it is not cancellation of vertices or of real heights.

If two points on one side coincide in $G_H$, discard the node as a
forced collision. If $B=s+\epsilon A$ in $G_H$, with
$\epsilon\in\{1,-1\}$, discard it as forced $\TI$. Both conditions
persist in every quotient. Otherwise an empty residual list is a
homometric terminal. At a nonterminal choose one residual $A$ edge
$(i,j)$. For each residual $B$ class use its first labelled edge
$(k,l)$ and include each sign $\epsilon$ satisfying
\[
 h_j-h_i=\epsilon(h_l-h_k).
\]
The child adjoins the actual integral row
$\alpha_{ij}^A-\epsilon\alpha_{kl}^B$. Both signs are included if
both satisfy this equality, as happens for zero heights.

\begin{lemma}[Coverage and termination of completion]
Suppose a full presentation $G_M$ in this stratum is a quotient of a
node $G_H$. Unless its realization has a collision or is $\TI$,
it factors through one recorded child of every nonterminal encountered.
Every path terminates after at most fifteen uncancelled-edge decreases.
\end{lemma}
\begin{proof}
The cancelled classes remain equal in $G_M$. Subtracting the same
multiset from the two equal full edge multisets therefore leaves equal
residual multisets in $G_M$. The chosen $A$ occurrence equals some
residual $B$ occurrence up to sign there. Replacing that occurrence by
the saved representative of its unsigned class changes at most the
sign, since equality up to sign already holds in $G_H$. Both possible
signs are tested. The common real height homomorphism is defined on
$G_M$, so the resulting signed height equality is among the branches.
The newly adjoined row holds in $G_M$ and gives the required factorization.

Existing unsigned equalities persist on passing to a quotient. The new
equality pairs one occurrence from each residual side that had not been
paired before. Even if further classes merge, a matching of occurrences
extending all previous cancellations and this new pair exists. Hence the
maximum cancellable multiplicity increases by at least one, and the
residual multiplicity strictly decreases. It starts at most fifteen.
Two-way lattice certificates make cache reuse a reuse of the same
labelled quotient, so this decrease also precludes cycles after merging.
\end{proof}

This argument uses the original presentation and its real homomorphism;
it does not posit a nonzero homomorphism from the finite cyclic target
back to $\R$. The root-orbit reduction may first replace $M$ by its
within-fibre relabelling; that presentation has the same realization
up to the recorded changes of labels and anchors.

Here is a complete abstract version of the completion checker:
\begin{verbatim}
for each height stratum:
  verify height orbit cover, generic heights, and cross-root orbit cover
  for each root: check both integral inclusions into its saved node
  visit each reachable node once:
    verify Smith coordinates and all point/edge classes
    recompute residual lists by paired multiplicity cancellation
    if collision: verify two distinct labels have equal coordinates
    elif T/I: verify the saved sign and shift give the entire other set
    elif homometric: require empty residuals and six points on both sides
    else:
      require chosen edge belongs to the residual A list
      regenerate every height-compatible (B representative, sign)
      require the saved branch list equals this list, without duplicates
      for each branch: check its exact row and both lattice inclusions
      require the child's residual multiplicity is strictly smaller
  require every saved node is reachable and every root orbit is represented
\end{verbatim}
The package must report \code{complete=true}, no error, and certified
normalizations; a capped or interrupted search is rejected. The accepted
cover has $10\,602$ nodes: $7\,429$ forced $\TI$, $2\,207$ collisions,
$620$ homometric terminals and $346$ branching nodes, with $974$ branch
edges. The $658\,894$ integral inclusion identities are exactly twice
the number of roots plus branch edges.

\paragraph{Terminal mechanism predicates.}
For every homometric terminal the mechanism atlas must contain exactly
one record with the same stratum, node identifier, points and group orders.
Write $\widetilde B=s+\epsilon B$ for its certified alignment. For a
block certificate check a nonempty disjoint partition $A=U\sqcup W$,
its proposed moved block $W'$, and $\widetilde B=U\sqcup W'$.
With $P=UW^*$, the exact integer coefficient tests are
\[
\begin{array}{c|c|c}
\text{label}&W'&\text{acceptance identity}\\\hline
\mathrm{L2}&x^tW&x^tP=P\\
\mathrm{L3*}&x^tW&x^{-t}P=P^*\\
\mathrm{L4}&x^tW&2t=0,\quad x^t(P+P^*)=P+P^*\\
\mathrm{L5}&x^tW^*&P=x^{-t}UW .
\end{array}
\]
The group-ring soundness proofs given above apply literally. A verifier
expands these distributions as integer counters; it never accepts a label
merely because an endpoint autocorrelation check succeeds.

A Bloom certificate supplies $p,q$, anchor and alignment. Evaluate the
six displayed Bloom formulas and require exact endpoint equality with
multiplicity. A D certificate supplies an anchor, four-set $C$ and steps
$a,b$. Check the two endpoint formulas and expand
\[
 K=C+x^{a+b}C^*+(1+x^a)(1+x^b),\qquad
 (1-x^a)(1-x^b)K=0.
\]
The cyclic periodic-kernel lemma proved above converts this exact identity
to the stated nonnegative periodic-weight D grammar. If specialization
makes the two equal-sum dyads share a point, they share both points and
the endpoints agree; otherwise the four corners remain distinct.

An L7 certificate supplies a generator $h$ of a subgroup $H$ of order
$m\in\{2,4,6,12\}$. Enumerate its $m$ multiples and check their
distinctness and $mh=0$. Form $S=A+H$, require $|S|=12$, check exactly
$m/2$ points of $A$ in every occupied $H$-coset, and require
$\widetilde B=S\setminus A$. Its previously proved specialization
lemma is essential: a surviving non-$\TI$ six-point cyclic image
preserves the subgroup order and occupied cosets. Thus the image is a
literal L7 construction, not only a transported polynomial identity.

The producer's accepted $611$ direct terminal certificates comprise
$271$ L4, $135$ L3*, $114$ L7, $53$ L5, $16$ L2, $14$ Bloom and $8$ D.
A different independent search priority supplies $481$ D, $71$ L5,
$24$ L3*, $20$ L4, $14$ Bloom and one L2, again leaving exactly nine
terminals. These distributions are not invariants or a disjoint
classification; acceptance depends on the supplied identities.

\paragraph{All cyclic images of noncyclic torsion.}
For any remaining terminal write $G_H=T\oplus\Z^r$ with
$T=\bigoplus_{i=1}^k C_{d_i}$ and $E=\operatorname{lcm}(d_i)$.
Enumerate all tuples $j_i\in\{0,\ldots,d_i-1\}$, set
$c_i=(E/d_i)j_i$, $g=\gcd(E,c_1,\ldots,c_k)$ and $q=E/g$.
The character has actual image order $q$, and its surjective image map is
\[
 \theta_j(t_1,\ldots,t_k)=
 \frac{\sum_i c_it_i}{g}\pmod q.
\]
Retain the free coordinates unchanged. If $q=1$ omit the torsion coordinate.

\begin{lemma}[Actual-image character factorization]
Every map $f:T\oplus\Z^r\to C_n$ factors through one of the displayed
quotients $C_q\oplus\Z^r$, with $q\mid n$.
\end{lemma}
\begin{proof}
The image $f(T)$ is cyclic of order $q$ dividing both $E$ and $n$.
Identify it with $C_q$, and embed $C_q$ in $C_E$ by multiplying by
$E/q$. Each generator of order $d_i$ then maps to an element killed by
$d_i$, hence to $(E/d_i)j_i$ for one enumerated $j_i$. Dividing the
resulting character by its common divisor $g=E/q$ recovers the chosen
map onto $C_q$. Extend its factor map by the original images of the free
generators. These images are unrestricted and need not be disjoint from
the torsion image. Requiring $C_E$ itself to embed into $C_n$ would be
incorrect; only the actual image $C_q$ is used.
\end{proof}

The checker requires every character tuple exactly once, recomputes
$c_i,E,g,q$ and all projected point coordinates, then verifies a collision,
a $\TI$ witness or one of the preceding mechanism certificates.
Seven residual terminals have $T=C_2\oplus C_2$: each of their four
characters forces a collision. Two have $T=C_2\oplus C_8$: each has
sixteen characters, eight colliding and eight admitting L7 certificates
(independently also D). Thus all $60$ quotients give $44$ collisions and
$16$ mechanisms. Coverage, specialization and this factorization prove LR.

\subsection{The bounded finite-source branch}
\begin{lemma}[FT]\label{lem:ft}
For every $6\le q\le135$, each homometry family of six-subsets of
$C_q$ is connected by the stated grammar.
\end{lemma}
\begin{proof}[Proof with its finite acceptance contract]
Enumerate six-subsets in two independent ways. First enumerate every
positive gap word $(g_0,\ldots,g_5)$ with sum $q$, keeping precisely the
lexicographically least word among its rotations and reversals. Equivalently
require $g_0$ minimal and then test all twelve gap images. Cumulative sums
give its anchored point tuple. Gap rotations and reversals are exactly
translations and reflections of the cyclic point set, so this visits every
$\TI$ class once, including symmetric classes. The second traversal
enumerates every $0<a_1<\cdots<a_5<q$ and keeps the least sorted tuple
among all images anchored at one of its six points, with either sign.
Any rigid image containing zero has such an anchor, so this test is complete.
Lexicographic comparison of point tuples and their gap words first differs
at the same gap, so the two canonical conventions agree.

For the first traversal compute the fifteen distances
$\min(a_j-a_i,q-a_j+a_i)$. For the second compute directed membership
counts $\#\{x\in A:x+d\in A\}$ for $1\le d\le q/2$, halving the
antipodal entry. Expand these counts into the sorted fifteen-distance list.
Both lists are exact interval inventories. Group complete inventories,
retaining groups of at least two classes. The implementation packs fifteen
distance bytes in two words and five point bytes in another; for
$q\le135$ these encodings are injective. Equality uses both full distance
words, not a hash. Two traversals must agree family-for-family and in
total class count at every modulus, not merely in aggregate totals.

An additional class-count check uses Burnside's formula. A rotation by
$t$ has $c=\gcd(q,t)$ cycles of length $\ell=q/c$ and fixes
$\binom{c}{6/\ell}$ six-subsets when $\ell\mid6$, otherwise zero.
A reflection with $f$ fixed points fixes
$\sum_{j:(6-j)\text{ even}}\binom f j
\binom{(q-f)/2}{(6-j)/2}$, where $f=1$ for odd $q$, and $f=2$ or
$0$ according to even or odd reflection parameter for even $q$.
Sum over all $2q$ group elements and divide by $2q$.

Each census checkpoint records \code{methods}, \code{summary},
\code{independent}, \code{families} and \code{source_sha256}.
Each family stores its complete \code{members} and \code{icv}.
Reference regressions check the immutable implementation through $22$;
additional controls attack the $31$--$33$ and $63$--$65$ boundaries.

For every unordered pair within every family, the mechanism file stores
either a direct certificate or an explicit gap entry. Certified pairs and
gaps must partition the complete pair set, without duplicates or omissions.
A direct record stores canonical \code{a}, \code{b} and \code{label}.
Bloom records store \code{parameters} $[p,q']$; R records store a seed
index and multiplier, with its seed-order relation checked. L7 records
store subgroup order \code{h}; unit records store \code{u}, with
$\gcd(u,q)=1$, endpoint equality and invariant autocorrelation checked.
Block/D records store \code{certificate} $[\lambda,\epsilon,s,\mu,t]$:
$\mu$ is a six-bit mask selecting $U$ from the recorded $A$, while
$\epsilon,s$ align $B$. Reconstruct both blocks and verify the preceding
coefficient identities. The older \code{swap} label additionally requires
its halved-shift solvability condition; it is therefore contained in L3*.
For D reconstruct its equal-sum dyads and check the displayed mixed
difference of its twelve-unit weight distribution. The cyclic kernel lemma
then supplies the required nonnegative periodic decomposition.

Finally, form an undirected graph on each family's members with its
verified direct certificates as edges. Require one connected component,
using an independent union-find or exhaustive reachability calculation.
For every gap pair save and verify an explicit path of these edges; no
direct construction is required for that pair. The accepted tables have
$725\,132$ families and $728\,424$ unordered pair edges, $728\,351$
direct certificates and $73$ gaps, all connected by verified paths.
Moduli $6$--$11$ have no nontrivial family. These exact finite conditions,
together with soundness of the generators, prove FT.
\end{proof}

\begin{lemma}[Inflation]\label{lem:inflate}
If $q\mid n$, any generating path of six-subsets of $C_q$ embeds to
a generating path in $C_n$ under $x\mapsto(n/q)x$. Distinct $\TI$
classes stay distinct. Consequently FT and the rank-ten finite-source
bound settle the original free-rank-zero branch at every modulus.
\end{lemma}
\begin{proof}
The embedding is injective. It preserves every point distinction,
partition, subgroup, periodic-weight distribution and group-ring identity,
so it transports Bloom, D, block, coset and rigid-template edges. For a
unit $u$ modulo $q$, choose $u'$ congruent to $u$ modulo $q$ and congruent
to $1$ modulo every prime dividing $n$ but not $q$. The moduli in these
extra congruences are coprime to $q$, so CRT supplies $u'$. Primes dividing
$q$ cannot divide $u'$ because they do not divide $u$; the other primes
were excluded explicitly. Thus $u'$ is a unit modulo $n$ and acts on
the embedded subgroup exactly as $u$. Autocorrelation is zero off that
subgroup, hence the multiplier condition also transports.

If two nonempty subsets of the embedded subgroup $H$ satisfy
$B=s+\epsilon A$ in $C_n$, choose one point on each side in this equality.
Their difference shows $s\in H$. The same equality is therefore a rigid
equivalence inside $H\cong C_q$, proving preservation of vertices.
For an original rank-ten presentation $G_M$ the group is finite, with
$|G_M|\le135$ by the minor bound. Its cyclic image has order
$6\le q\le135$ and divides $n$. Independently anchored inputs lie
inside this subgroup, so FT applies and its path embeds as just proved.
Restore the original anchor translations at the endpoints.
\end{proof}

\paragraph{Reproducibility and review boundary.}
The height tables are in \code{results/2026-09-30-six-free-rank/};
the compressed cross/root/node certificates are in
\code{results/2026-09-30-six-free-dag-all/}; the terminal atlas is
\code{results/2026-09-30-six-free-mechanisms/atlas.json}.
The structural checker is \code{tests/test_six_free_dag_review.py};
the independent terminal checker is
\code{tests/test_six_low_rank_mechanisms_review.py}.
They check different obligations and both are required. The additional
height checker is \code{tests/test_six_free_rank_growth.py}.
Finite census and edge checkpoints are respectively in
\code{results/2026-09-30-six-large-census/} and
\code{results/2026-09-30-six-pair-mechanisms/}; the finite bridge audit is
\code{tests/test_six_finite_torsion_review.py}.
The supplied independent review of the earlier draft rechecked height
counts and consequences but explicitly did not replay the internal LR DAG.
Its review must not be cited as such a replay. The separately implemented
in-house structural and terminal checks are the certificate evidence for
that branch. These finite computational claims remain part of a
computer-assisted proof, not a proof-assistant-kernel or novelty claim.
